\documentclass[../main.tex]{subfiles}
\begin{document}

\section*{Results}
The project delivered a working project management system that meets its objectives:
\begin{itemize}
  \item Accounts without global roles; projects with per-project PM and Member roles, invitations
        that must be accepted, and authorisation enforced by the server on every route.
  \item A keyboard-first Kanban board and table view with deadline-first ordering, several
        assignees per task, dependencies, a critical path and a DAG view.
  \item The Scrum process: sprint lifecycle, backlog planning, burndown, velocity and an AI
        retrospective; plus a change history on every task.
  \item Five AI features backed by Groq with a three-tier fallback, structured and validated output,
        and answers in the interface language.
  \item Local-first behaviour with a safe offline policy; a complete English/Vietnamese interface.
  \item A production deployment on \url{https://koshi.huynhca.com}, 185 backend and 442 frontend
        tests, continuous integration, and eight living documents that trace every requirement to
        code and tests.
  \item An AI agent ``software factory'': six specialised agent stations, guardrails that enforce
        the project's policies on every edit and command (57 tests), and health checks, request
        ids and metrics for operating the product.
\end{itemize}

Beyond the features, the most valuable lesson concerned verification. Several of the most serious
defects---a secret shipped inside the container image, a deploy command that would have overwritten
production with development settings, an ``AI'' feature that never called a model, a language
switch that left half the interface untranslated---were invisible to a test suite that was green.
They were found by checking the real artefact: inspecting the built image, reading the command
before running it, calling the real model, and scanning every screen. The practice of
mutation-testing the tests, and of recording every decision with its reason, came from these
experiences.

\section*{Limitations}
\begin{itemize}
  \item Shared projects are read-only offline; there is no algorithm for merging concurrent offline
        edits.
  \item The burndown is computed from a sprint's current membership, so work moved out of a sprint
        disappears from its history.
  \item The deterministic AI fallback answers only in Vietnamese.
  \item The public instance has no login rate limiting, and the performance targets in the
        requirements (rendering within one frame, memory use) have not been measured.
  \item The public instance runs on a single machine and is unavailable when that machine is off.
\end{itemize}

\section*{Future work}
\begin{itemize}
  \item Conflict-free replicated data types (CRDTs) so shared projects can also be edited offline.
  \item An event log per sprint for exact burndown history, and a daily-scrum view.
  \item Rate limiting and account lockout; managed hosting with automatic backups.
  \item Native mobile clients and integration with Git hosting (linking commits and pull requests
        to tasks automatically).
\end{itemize}

\end{document}
